% Generated draft: requires mathematical and source review.
\documentclass[11pt]{article}
\usepackage[margin=25mm]{geometry}
\usepackage{fontspec}
\setmainfont{Latin Modern Roman}
\usepackage{amsmath,amssymb,mathtools,unicode-math}
\setmathfont{Latin Modern Math}
\usepackage{enumitem,needspace,xurl,hyperref}
\setlength{\parindent}{0pt}
\setlength{\parskip}{5pt}
\setlength{\emergencystretch}{3em}
\newcommand{\N}{\mathbb{N}}
\newcommand{\Z}{\mathbb{Z}}
\newcommand{\Q}{\mathbb{Q}}
\newcommand{\R}{\mathbb{R}}
\newcommand{\C}{\mathbb{C}}
\newcommand{\F}{\mathbb{F}}
\newcommand{\A}{\mathbb{A}}
\newcommand{\PP}{\mathbb{P}}
\newcommand{\E}{\mathbb{E}}
\newcommand{\eps}{\varepsilon}
\newcommand{\dd}{\,\mathrm{d}}
\begin{document}
\section*{UnsolvedMath ID: 30005502}
\textbf{Source problem code:} catalogue entry
\subsection{Definitions and mathematical statement}
Let \(p\) be a prime, let \(\F_p\) be the field with \(p\) elements, and let \(S_n\) be the symmetric group on \(n\) letters. The group algebra \(A=\F_pS_n\) is the \(\F_p\)-vector space with basis \(S_n\), with multiplication obtained by extending the group multiplication bilinearly.

For a unital associative algebra \(A\) over a field \(k\), let
\[
\operatorname{Der}_k(A)=\{D\in\operatorname{End}_k(A):D(ab)=D(a)b+aD(b)\text{ for all }a,b\in A\}.
\]
Let \(\operatorname{Inn}(A)\) be the subspace of derivations of the form \(a\mapsto xa-ax\), for \(x\in A\). The first Hochschild cohomology is
\[
HH^1(A)=\operatorname{Der}_k(A)/\operatorname{Inn}(A),
\]
with Lie bracket induced by \([D,E]=D\circ E-E\circ D\). For a Lie algebra \(L\) over \(k\), define its derived series by
\[
L^{(0)}=L,\qquad L^{(j+1)}=\operatorname{span}_k\{[x,y]:x,y\in L^{(j)}\}.
\]
The Lie algebra is solvable if \(L^{(d)}=0\) for some integer \(d\geq0\).

\textbf{Solvability conjecture.} For every prime \(p\) and every integer \(n\geq1\), the Lie algebra \(HH^1(\F_pS_n)\) is solvable. Equivalently,
\[
\forall p\text{ prime}\;\forall n\in\Z_{\geq1}\;\exists d\in\Z_{\geq0},
\qquad \bigl(HH^1(\F_pS_n)\bigr)^{(d)}=0.
\]
The integer \(d\) may depend on both \(p\) and \(n\); no uniform bound is asserted.

\textbf{Associated generating-function question.} For an algebra \(A\) over \(k\), higher Hochschild cohomology is \(HH^r(A)=\operatorname{Ext}_{A\otimes_k A^{\mathrm{op}}}^{r}(A,A)\), where \(r\geq0\). A block algebra of a finite-dimensional algebra is a direct summand \(Ae\) associated with a primitive central idempotent \(e\). The source asks whether the preceding results on generating functions for Hochschild cohomology can be extended to \emph{(blocks of) alternating groups, other finite Coxeter groups, Hecke algebras, and diagram algebras}. For finite groups, the algebraic setting is that of their group algebras. This is an exploratory question, rather than a specified generating-function identity.

The supplied extraction includes a damaged formula apparently involving coefficients \(\dim_{\F_p}HH^r(kS_n)\), powers \(t^n\), a sum beginning at \(n=0\), and an expression \(t^p\phi(t^p)P(t)\). It does not define \(k\), the intended range or role of \(r\), \(\phi\), or \(P\). The preceding generating-function results are also absent. Accordingly, no equality is asserted from this fragment. The algebra families, parameters, coefficient fields, block choices, and intended identities for the extension question require clarification from the source.
\subsection{Short English statement}
The conjecture asks whether, for every prime \(p\) and every positive integer \(n\), repeated derived Lie brackets eventually vanish in the space of derivations of \(\F_pS_n\) modulo inner derivations. The source also asks for extensions of Hochschild-cohomology generating functions to further algebra families and blocks, but the supplied record does not specify enough data to formulate those extensions as precise identities.
\subsection{Sources}
\begin{itemize}
\item Local UnsolvedMath registry, \texttt{data/unsolvedmath/problems.json}, permanent ID \(30005502\), title: \emph{Hochschild Cohomology of Symmetric-Group Algebras}; exact source problem code: \texttt{catalogue entry}.
\item Supplied report citation: \emph{Representations of Finite Groups} (2023), Oberwolfach Reports, DOI 10.4171/owr/2023/19, pp.\ 1031--1084. \url{https://doi.org/10.4171/owr/2023/19}. The original extraction attributes the conjecture to ``the authors'' without naming them; the registry gives proposed year 2023 and no named proposer.
\item Supplied literature reference: Linckelmann--Wang, \emph{solvability of HH1 for blocks}. \url{https://arxiv.org/abs/2503.16697}.
\item Supplied literature reference: Benson--Kessar--Linckelmann, \emph{Hochschild cohomology II}. \url{https://arxiv.org/abs/2301.03185}.
\item Supplied literature reference: \emph{Oberwolfach Report 19/2023}. \url{https://doi.org/10.4171/owr/2023/19}.
\item Literature assessment dated \texttt{2026-08-22}: the supplied curation labels the problem open and reports substantial criteria and computations, including recent block-theoretic criteria, but no established general proof for every \(n\) and \(p\). This is supplied dated curation derived from the Oberwolfach record, not an independently verified literature assessment. The record also labels its statement \texttt{corrected\_verified} and says it was reconstructed and checked against the official report; those are supplied verification claims, not verification performed for this draft.
\end{itemize}
\end{document}
